% =========================================================
% PART III -- SAMPLE STATISTICS
% =========================================================

\section{Sample Statistics}

% ---------------------------------------------------------
% 14. POPULATION AND SAMPLE
% ---------------------------------------------------------

\begin{frame}{14. Population and Sample}

\textbf{Population}

A population is the \textbf{complete set of individuals,
objects, or observations} of interest in a study.

\vspace{0.4cm}

\textbf{Sample}

A sample is a \textbf{subset of the population} selected for
analysis.

\vspace{0.4cm}

\[
\boxed{
\text{Population}
\supseteq
\text{Sample}
}
\]

\vspace{0.3cm}

\textbf{Example}

If a university has 5,000 students and we select 200 students
for analysis:

\begin{itemize}
    \item Population = 5,000 students
    \item Sample = 200 selected students
\end{itemize}

\end{frame}

% ---------------------------------------------------------
% 14.1 POPULATION AND SAMPLE -- FIGURE
% ---------------------------------------------------------

\begin{frame}{14.1 Population and Sample -- Figure}

\begin{columns}[T]

\begin{column}{0.55\textwidth}

\centering
\includegraphics[width=\textwidth]{figures/population_sample.png}

\end{column}

\begin{column}{0.46\textwidth}

The figure illustrates the relationship between a population, a sampling frame, and a sample in a research study.

\vspace{0.1cm}

The large gray circle represents the population, the entire target population under study.

\vspace{0.1cm}

The blue circle inside it represents the sampling frame, the portion of the target population from which participants can be selected.

\vspace{0.1cm}

The dark navy circle represents the sample, the actual units or individuals selected for the study.

\end{column}

\end{columns}

\end{frame}

% ---------------------------------------------------------
% 14.2 POPULATION VS SAMPLE
% ---------------------------------------------------------

\begin{frame}{14.2 Population and Sample -- Comparison}

\small

\begin{table}
\centering

\begin{tabular}{>{\bfseries}l l l}
\toprule
 & Population & Sample \\
\midrule
Size & Complete set & Subset \\
Mean & $\mu$ & $\bar{x}$ \\
Variance & $\sigma^2$ & $s^2$ \\
Standard deviation & $\sigma$ & $s$ \\
Purpose & Describe population & Estimate population \\
\bottomrule
\end{tabular}

\end{table}

\vspace{0.5cm}

\textbf{Key idea}

Since studying an entire population may be expensive or
impractical, statistical analysis often uses a sample to
learn about the population.

\vspace{0.3cm}

\textbf{Application}

Sampling is fundamental in machine learning, surveys,
experiments and data analysis.

\end{frame}


% ---------------------------------------------------------
% 15. SAMPLE MEAN
% ---------------------------------------------------------

\begin{frame}{15. Sample Mean}

\textbf{Definition}

The sample mean is the arithmetic average of the observations
in a sample.

For a sample
\[
x_1,x_2,\ldots,x_n
\]

the sample mean is
\[
\boxed{
\bar{x}
=
\frac{1}{n}
\sum_{i=1}^{n}x_i
}
\]

\vspace{0.1cm}

\textbf{Interpretation}

The sample mean provides an estimate of the population mean $\mu$.

\[
\boxed{
\bar{x}\approx\mu
}
\]

The quality of this estimate generally improves as the sample
size increases.

\end{frame}


% ---------------------------------------------------------
% 15.1 SAMPLE MEAN -- EXAMPLE
% ---------------------------------------------------------

\begin{frame}{15.1 Sample Mean -- Example}

\textbf{Example}

Consider a sample of five observations:

\[
X=\{12,15,18,20,25\}
\]

The sample mean is

\[
\bar{x}
=
\frac{12+15+18+20+25}{5}
\]

\[
\boxed{
\bar{x}=18
}
\]

\vspace{0.4cm}

\textbf{Interpretation}

The value $18$ is the average of the observations in the
sample and can be used as an estimate of the population mean.

\end{frame}


% ---------------------------------------------------------
% 16. SAMPLE VARIANCE
% ---------------------------------------------------------

\begin{frame}{16. Sample Variance}

\textbf{Definition}

Sample variance measures the variability of observations
around the sample mean.

\vspace{0.4cm}

For a sample $x_1,x_2,\ldots,x_n$:

\[
\boxed{
s^2
=
\frac{1}{n-1}
\sum_{i=1}^{n}
(x_i-\bar{x})^2
}
\]

\vspace{0.4cm}

\textbf{Why $n-1$?}

The denominator $n-1$ provides an unbiased estimator of the
population variance under the usual random-sampling assumptions.

\vspace{0.3cm}

\textbf{Key idea}

Sample variance measures how much the observations vary
around their sample mean.

\end{frame}


% ---------------------------------------------------------
% 16.1 SAMPLE VARIANCE -- EXAMPLE
% ---------------------------------------------------------

\begin{frame}{16.1 Sample Variance -- Example}

\textbf{Example}

Consider:
\[
X=\{2,4,6\}
\]
The sample mean is
\[
\bar{x}=4
\]
Squared deviations:
\[
(2-4)^2=4,\qquad
(4-4)^2=0,\qquad
(6-4)^2=4
\]
\[
s^2
=
\frac{4+0+4}{3-1}
\]
\[
\boxed{s^2=4}
\]
\textbf{Interpretation}

The sample variance quantifies the spread of the observations
around their sample mean.

\end{frame}


% ---------------------------------------------------------
% 17. ORDER STATISTICS
% ---------------------------------------------------------

\begin{frame}{17. Order Statistics}

\textbf{Definition}

Order statistics are the observations arranged in
\textbf{ascending order}.

\vspace{0.1cm}

Given observations
\[
x_1,x_2,\ldots,x_n,
\]

their ordered values are written as
\[
\boxed{
x_{(1)}\leq x_{(2)}\leq\cdots\leq x_{(n)}
}
\]

where $x_{(i)}$ denotes the $i$-th smallest observation.

\vspace{0.1cm}

\textbf{Examples of order statistics}

\begin{itemize}
    \item $x_{(1)}$ -- minimum
    \item $x_{(n)}$ -- maximum
    \item Middle order statistic -- median
    \item Quartiles -- divide ordered data into parts
\end{itemize}

\end{frame}


% ---------------------------------------------------------
% 17.1 ORDER STATISTICS -- EXAMPLE
% ---------------------------------------------------------

\begin{frame}{17.1 Order Statistics -- Example}

\textbf{Example}

Consider:
\[
X=\{40,10,30,50,20\}
\]

Arrange the observations in ascending order:
\[
10,\;20,\;30,\;40,\;50
\]
Therefore,
\[
x_{(1)}=10,\qquad
x_{(2)}=20,\qquad
x_{(3)}=30,\qquad
x_{(4)}=40,\qquad
x_{(5)}=50
\]
Hence,
\[
\boxed{x_{(1)}=\min(X)}
\]
\[
\boxed{x_{(5)}=\max(X)}
\]

The middle observation $x_{(3)}=30$ is the median.

\end{frame}


% ---------------------------------------------------------
% 18. COVARIANCE
% ---------------------------------------------------------

\begin{frame}{18. Covariance}

\textbf{Definition}

Covariance measures the \textbf{direction of the linear
relationship} between two variables.

\vspace{0.4cm}

For random variables $X$ and $Y$:

\[
\boxed{
\operatorname{Cov}(X,Y)
=
E[(X-\mu_X)(Y-\mu_Y)]
}
\]

\vspace{0.4cm}

\textbf{Interpretation}

\begin{itemize}
    \item Positive covariance: variables tend to increase together.
    \item Negative covariance: one tends to increase when the other decreases.
    \item Covariance near zero: weak linear association.
\end{itemize}

\vspace{0.3cm}

\textbf{Note}

The magnitude of covariance depends on the units of the variables.

\end{frame}


% ---------------------------------------------------------
% 18.1 COVARIANCE -- EXAMPLE
% ---------------------------------------------------------

\begin{frame}{18.1 Covariance -- Example}

\textbf{Example}

Consider two variables:
\[
X=\{1,2,3\}
\]
\[
Y=\{2,4,6\}
\]
As $X$ increases, $Y$ also increases.
Their means are:
\[
\bar{x}=2,\qquad \bar{y}=4
\]
The deviations have the same sign:
\[
(-1)(-2),\quad (0)(0),\quad (1)(2)
\]
All products are non-negative.
Therefore, the covariance is \textbf{positive}.

\textbf{Interpretation}

The variables show a positive linear relationship.

\end{frame}


% ---------------------------------------------------------
% 19. SAMPLE COVARIANCE
% ---------------------------------------------------------

\begin{frame}{19. Sample Covariance}

\textbf{Definition}

Sample covariance measures the linear relationship between
two variables using sample observations.

\vspace{0.1cm}

For paired observations
\[
(x_1,y_1),(x_2,y_2),\ldots,(x_n,y_n),
\]

the sample covariance is
\[
\boxed{
s_{XY}
=
\frac{1}{n-1}
\sum_{i=1}^{n}
(x_i-\bar{x})(y_i-\bar{y})
}
\]

\vspace{0.1cm}

\textbf{Interpretation}

\begin{itemize}
    \item $s_{XY}>0$: positive association.
    \item $s_{XY}<0$: negative association.
    \item $s_{XY}\approx0$: weak linear association.
\end{itemize}

\end{frame}

% ---------------------------------------------------------
% 19.1 SAMPLE COVARIANCE -- EXAMPLE
% ---------------------------------------------------------

\begin{frame}{19.1 Sample Covariance -- Example}

\textbf{Example}

Consider two variables:
\[
X=\{1,2,3\},
\qquad
Y=\{2,4,6\}
\]

The sample means are:
\[
\bar{x}=2,
\qquad
\bar{y}=4
\]

\begin{table}
\centering
\small

\begin{tabular}{c c c c c}
\toprule
$x_i$ & $y_i$ & $x_i-\bar{x}$ & $y_i-\bar{y}$ &
Product \\
\midrule
1 & 2 & $-1$ & $-2$ & $2$ \\
2 & 4 & $0$ & $0$ & $0$ \\
3 & 6 & $1$ & $2$ & $2$ \\
\bottomrule
\end{tabular}

\end{table}
\[
s_{XY}
=
\frac{2+0+2}{3-1}
=
\boxed{2}
\]

Since $s_{XY}>0$, $X$ and $Y$ have a positive covariance
and tend to increase together.

\end{frame}

% ---------------------------------------------------------
% 20. CORRELATION
% ---------------------------------------------------------

\begin{frame}{20. Correlation}

\textbf{Definition}

Correlation is a standardized measure of the
\textbf{strength and direction of linear association}
between two variables.

\vspace{0.1cm}

The population correlation coefficient is
\[
\boxed{
\rho_{XY}
=
\frac{\operatorname{Cov}(X,Y)}
{\sigma_X\sigma_Y}
}
\]

For a sample, the corresponding measure is the
\textbf{sample correlation coefficient} $r$.
\[
\boxed{
-1\leq r\leq1
}
\]

\textbf{Interpretation}

\begin{itemize}
    \item $r=1$: perfect positive linear relationship.
    \item $r=-1$: perfect negative linear relationship.
    \item $r\approx0$: weak or no linear relationship.
\end{itemize}

\end{frame}


% ---------------------------------------------------------
% 20.1 CORRELATION -- COVARIANCE
% ---------------------------------------------------------

\begin{frame}{20.1 Correlation -- Relation to Covariance}

The sample correlation coefficient can be written as
\[
\boxed{
r
=
\frac{s_{XY}}{s_Xs_Y}
}
\]
where:

\begin{itemize}
    \item $s_{XY}$ = sample covariance
    \item $s_X$ = sample standard deviation of $X$
    \item $s_Y$ = sample standard deviation of $Y$
\end{itemize}

\vspace{0.2cm}

\textbf{Key difference}

\begin{itemize}
    \item Covariance depends on the units of measurement.
    \item Correlation is standardized and always lies between $-1$ and $1$.
\end{itemize}

\vspace{0.1cm}

\textbf{Application}

Correlation is commonly used for feature analysis and
understanding relationships between variables in datasets.

\end{frame}


% ---------------------------------------------------------
% 21. SAMPLE COVARIANCE MATRIX
% ---------------------------------------------------------

\begin{frame}{21. Sample Covariance Matrix}

\textbf{Definition}

For a dataset containing multiple variables, the sample
covariance matrix summarizes the covariance between every
pair of variables.

For $p$ variables:
\[
\boxed{
S=
\begin{bmatrix}
s_{11} & s_{12} & \cdots & s_{1p}\\
s_{21} & s_{22} & \cdots & s_{2p}\\
\vdots & \vdots & \ddots & \vdots\\
s_{p1} & s_{p2} & \cdots & s_{pp}
\end{bmatrix}
}
\]

where
\[
s_{ij}
=
\frac{1}{n-1}
\sum_{k=1}^{n}
(x_{ki}-\bar{x}_i)(x_{kj}-\bar{x}_j)
\]

The diagonal elements represent the \textbf{sample variances}
of the variables.

\end{frame}


% ---------------------------------------------------------
% 21.1 SAMPLE COVARIANCE MATRIX -- EXAMPLE
% ---------------------------------------------------------

\begin{frame}{21.1 Sample Covariance Matrix -- Example}

Suppose a dataset contains three variables:

\[
\{\text{Age, Study Hours, Marks}\}
\]

A sample covariance matrix may be represented as

\[
S=
\begin{bmatrix}
s_{\text{Age,Age}} &
s_{\text{Age,Study}} &
s_{\text{Age,Marks}}
\\
s_{\text{Study,Age}} &
s_{\text{Study,Study}} &
s_{\text{Study,Marks}}
\\
s_{\text{Marks,Age}} &
s_{\text{Marks,Study}} &
s_{\text{Marks,Marks}}
\end{bmatrix}
\]

\textbf{Interpretation}

\begin{itemize}
    \item Diagonal entries represent variances.
    \item Off-diagonal entries represent covariances.
    \item The matrix is symmetric:
    \[
    s_{ij}=s_{ji}
    \]
\end{itemize}

\end{frame}


% ---------------------------------------------------------
% 21.2 SAMPLE COVARIANCE MATRIX -- APPLICATION
% ---------------------------------------------------------

\begin{frame}{21.2 Sample Covariance Matrix -- Application}

\textbf{Why it is useful}

The covariance matrix provides a compact representation of
relationships among multiple numerical variables.

\vspace{0.4cm}

\textbf{Applications in AI and Computing}

\begin{itemize}
    \item Feature analysis
    \item Principal Component Analysis (PCA)
    \item Multivariate statistical analysis
    \item Understanding relationships between features
\end{itemize}

\vspace{0.4cm}

\textbf{Key idea}

\[
\boxed{
\text{Multiple variables}
\longrightarrow
\text{Covariance Matrix}
\longrightarrow
\text{Multivariate Analysis}
}
\]

\end{frame}